\chapter{Training Procedure}
\label{chap:train}

This chapter describes the procedure used to adapt pretrained embedding models for more efficient $A^*$ traversal. Section~\ref{sec:train_dataset} explains how successful graph traversals are converted into supervised preference examples. Sections~\ref{sec:train_objective} and~\ref{sec:train_matryoshka} introduce the ranking-based training objective and its extension to Matryoshka Representation Learning. Section~\ref{sec:train_validation} presents the search-based validation procedure, while Sections~\ref{sec:train_hyperparameters} and~\ref{sec:train_final_training} describe hyperparameter optimization and final model training. Finally, Section~\ref{sec:train_final_precomputation} explains how the trained embeddings are precomputed and stored for evaluation and inference.

\section{Supervision Construction}
\label{sec:train_dataset}

To fine-tune the embedding models, supervised training and validation datasets are constructed from the corresponding splits of the binary causal question dataset derived from MS MARCO by \citet{bluebaum:2024}. This construction is specific to model training and hyperparameter optimization, while the subsequent evaluation data remain unchanged. Each pretrained base model is used to generate its own supervision.

To construct a dataset for a given split, consider a causal graph $G=(V,E)$. Graph-node embeddings are first precomputed using the respective pretrained base model. For each cause--effect pair $(A,B)$ in the split, let $a,b\in V$ denote the corresponding graph nodes. An unbounded $A^*$ search is performed from \mbox{$a$ to $b$} using distance metric $d\in\{\mathrm{Cos},\mathrm{Euc}\}$, denoting cosine and Euclidean distance, respectively. If the search reaches the target $b$, the discovered path is retained. Let
\[
P=(p_0,\dots,p_k), \qquad p_0=a,\quad p_k=b,\quad k\geq1,
\]
denote such a path, where $p_i$ is the node at position $i$ and $k$ is the number of edges in the path. The $A^*$ path is used instead of a minimum-hop path obtained with BFS because minimizing the number of edges does not necessarily yield a path useful for embedding-guided search. In contrast, $A^*$ prioritizes nodes according to the embedding-space edge costs and heuristic estimates of the pretrained model. The resulting paths therefore provide supervision aligned with the search procedure that fine-tuning aims to improve.

For each distance metric $d$, let $\mathcal{D}_{\mathrm{train}}^d$ and $\mathcal{D}_{\mathrm{val}}^d$ denote the datasets constructed from the training and validation splits, respectively. Let $N^+(p_i)$ denote the set of outgoing successors of node $p_i$. For every step along a retained path, the path successor $p_{i+1}$ is treated as the preferred choice and compared with each alternative successor $n\in N^+(p_i)\setminus\{p_{i+1}\}$. For either split $s\in\{\mathrm{train},\mathrm{val}\}$, the corresponding quadruples are added as
\[
\mathcal{D}_s^d \gets \mathcal{D}_s^d \cup
\left\{
(p_i,p_k,p_{i+1},n)
\;\middle|\;
0\leq i<k,\;
n\in N^+(p_i)\setminus\{p_{i+1}\}
\right\}.
\]
Each quadruple is subsequently denoted by $(c,e,p,n)$, representing the current node, target node, preferred successor on the discovered path, and alternative successor, respectively. It defines a pairwise preference that encourages the model to assign a lower estimated traversal cost to the preferred successor than to the competing choice. Algorithm~\ref{alg:training-dataset} summarizes the construction procedure. The resulting datasets are cached after construction and reused throughout hyperparameter optimization and final model training, avoiding repeated $A^*$ searches for generating the same supervision.


\section{Ranking-Based Training Objective}
\label{sec:train_objective}

The ranking-based objective uses quadruples $(c,e,p,n)\in \mathcal{D}_{\mathrm{train}}^d$ to reshape the embedding space for graph traversal. The preferred successor $p$ should yield a lower cumulative distance to the target $e$ than the alternative successor $n$. The difference between the preferred and alternative traversal costs is defined as\looseness=-1
\[
\begin{aligned}
\Delta &= f(p)-f(n) \\
&= \bigl[g(p)+h(p)\bigr]-\bigl[g(n)+h(n)\bigr] \\
&=
\left[
d\bigl(\phi(c),\phi(p)\bigr)
+
d\bigl(\phi(p),\phi(e)\bigr)
\right]
-
\left[
d\bigl(\phi(c),\phi(n)\bigr)
+
d\bigl(\phi(n),\phi(e)\bigr)
\right].
\end{aligned}
\]
A value $\Delta<0$ indicates that the preferred traversal decision is favored. The ranking loss is obtained by applying a non-linear activation function $\sigma$ to $\Delta$
\[
\mathcal{L}_{\mathrm{rank}}=\sigma(\Delta).
\]
This work evaluates ReLU and GELU as activation functions. An additional regularization term encourages the embedding norms to remain close to one
\[
\mathcal{L}_{\mathrm{reg}}
=
\sum_{x\in\{c,e,p,n\}}
\left(\lVert\phi(x)\rVert_2-1\right)^2.
\]
For a single embedding dimension, the resulting loss is
\[
\mathcal{L}
=
\mathcal{L}_{\mathrm{rank}}+\mathcal{L}_{\mathrm{reg}}.
\]
The objective therefore optimizes embedding distances for causal graph traversal rather than semantic similarity.


\section{Matryoshka Representation Learning}
\label{sec:train_matryoshka}

To support efficient causal knowledge graph traversal at different embedding dimensions, the ranking objective is extended using Matryoshka Representation Learning~\cite{kusupati:2022}. Rather than optimizing only the full embedding vector, the objective is applied to multiple nested prefixes of each representation. Let $D_{\mathrm{model}}$ denote the native embedding dimensionality of the underlying model. The set of training dimensions is defined as
\[
\mathcal{M}
=
\{D_{\mathrm{model}}\}
\cup
\left\{
2^k
\;\middle|\;
k\in\mathbb{N},\;
2\leq 2^k<D_{\mathrm{model}}
\right\}.
\]
The dimension 768 is additionally included in $\mathcal{M}$ to enable direct comparison with MPNet and Granite, whose native dimensionality is 768~\cite{allmpnet:modelcard,granite:modelcard}.
\[
\mathcal{M}\gets\mathcal{M}\cup\{768\}.
\]
For each $m\in\mathcal{M}$, let
\[
\phi_m(x)=\phi(x)_{1:m}
\]
denote the first $m$ components of the embedding of concept $x$. Let $f_m$, $g_m$, and $h_m$ denote the corresponding $A^*$ evaluation, path-cost, and heuristic terms computed using $\phi_m$. The corresponding evaluation difference is
\[
\begin{aligned}
\Delta_m &= f_m(p)-f_m(n) \\
&= \bigl[g_m(p)+h_m(p)\bigr]-\bigl[g_m(n)+h_m(n)\bigr] \\
&=
\left[
d\bigl(\phi_m(c),\phi_m(p)\bigr)
+
d\bigl(\phi_m(p),\phi_m(e)\bigr)
\right]
\\
&\phantom{=}
-
\left[
d\bigl(\phi_m(c),\phi_m(n)\bigr)
+
d\bigl(\phi_m(n),\phi_m(e)\bigr)
\right].
\end{aligned}
\]
The dimension-specific ranking loss is
\[
\mathcal{L}_{\mathrm{rank}}^{(m)}
=
\sigma(\Delta_m).
\]
The overall training loss is
\[
\mathcal{L}_{\mathrm{total}}
=
\sum_{m\in\mathcal{M}}
\mathcal{L}_{\mathrm{rank}}^{(m)}
+
\mathcal{L}_{\mathrm{reg}}.
\]
Thus, every Matryoshka prefix contributes a ranking loss, while the norm regularization remains global. Applying the objective at multiple prefix lengths encourages information relevant to causal graph traversal to be concentrated in the earlier embedding dimensions. This enables a single fine-tuned model to support multiple embedding dimensions without additional training, allowing the representation size to be selected according to the desired trade-off between computational efficiency and heuristic quality.


\section{Validation Procedure}
\label{sec:train_validation}

Although the training objective optimizes local ranking decisions, validation measures whether the learned embedding space improves global graph traversal. At the beginning of each validation epoch, embeddings for all graph nodes are recomputed using the current state of the model and stored in a temporary lookup cache. For each distinct pair $(c,e)$ occurring in the supervision dataset $\mathcal{D}_{\mathrm{val}}^d$ within a validation batch, an $A^*$ search is performed from the current node $c$ to the target node $e$ using these updated embeddings and distance metric $d$. The number of nodes visited before reaching the target is recorded as $v(c,e)$. Let $P_{\mathrm{val}}^d$ denote the collection of evaluated pairs after batch-wise deduplication. The primary validation metric is the average number of visited nodes
$$
\text{AvgVisited}
=
\frac{1}{|P_{\mathrm{val}}^d|}
\sum_{(c,e)\in P_{\mathrm{val}}^d}
v(c,e).
$$
Lower values indicate that fewer graph nodes are visited on average before reaching the target and therefore correspond to more efficient traversal.


\section{Hyperparameter Optimization}
\label{sec:train_hyperparameters}

Before final training, model-specific hyperparameter optimization is performed with Optuna~\cite{akiba:2019}. The search jointly optimizes the learning rate, activation function (ReLU or GELU), and distance metric (cosine or Euclidean) using the Tree-structured Parzen Estimator sampler~\cite{bergstra:2011}. The learning rate is sampled logarithmically from $[2.5 \times 10^{-5}, 3 \times 10^{-5}]$. This interval was selected based on preliminary experiments indicating stable and competitive validation performance within this range. Each trial performs a short training run and is evaluated by the average number of nodes visited during validation $A^*$ searches. Early stopping terminates training when this metric no longer improves~\cite{prechelt:1998}, while Optuna's MedianPruner stops trials that underperform the median result of previously completed trials at the same training step~\cite{optuna:medianpruner}. For each embedding model, hyperparameter optimization first runs 30 completed or pruned trials and then continues one trial at a time. The search stops after five consecutive trials without a lower average number of visited nodes on validation. An improvement resets this counter, allowing up to 50 trials. Each trial runs for at most 10 epochs with an early-stopping patience of 3 validation epochs. Training uses the AdamW optimizer~\cite{loshchilov:2019} and a \mbox{ReduceLROnPlateau} scheduler~\cite{pytorch:reducelronplateau}, which halves the learning rate after one validation epoch without improvement. The effective batch size is fixed at 128, using gradient accumulation for smaller physical batches~\cite{ott:2018,lightning:gradientaccumulation}.

\section{Final Model Training}
\label{sec:train_final_training}

After hyperparameter optimization, the configuration with the lowest average number of visited nodes is selected for each embedding model. Final training uses the selected learning rate, activation function, and distance metric for up to 50 epochs with early stopping and a patience of 10 validation epochs. The checkpoint with the lowest average number of visited nodes on the validation set is retained as the final model and saved together with its training metadata.

\section{Final Embedding Precomputation}
\label{sec:train_final_precomputation}

After training, embeddings are precomputed for all graph-node concepts required during evaluation and inference. The node concepts are stored in a fixed order in a shared JSONL file. All embedding matrices follow this same ordering, allowing each node concept to be mapped to its row index and its embedding to be retrieved directly from the corresponding memory-mapped NumPy matrix~\cite{harris:2020}. The model first produces embeddings at its native dimensionality. For every Matryoshka dimension used in the experiments, the corresponding prefix of each full embedding is stored in a separate matrix. During graph traversal, $A^*$ loads the matrix for the selected model and dimension instead of computing embeddings with the encoder at runtime. This reduces inference time and enables direct indexed access to all node embeddings. The approach therefore trades additional disk space for lower runtime.
